\newpage\renewcommand{\abstractname}{Notations}
\addcontentsline{toc}{chapter}{Notations}
\begin{abstract}


    \begin{bluebox}
        Le présent rapport aborde principalement la prédiction de séries temporelles par apprentissage sur des jeux de données physiques.
        Elle utilise à ce titre de nombreux éléments mathématiques différents dont la notation est rappelée ci-dessous pour permettre une lecture plus fluide par la suite.
    \end{bluebox}

    \begin{itemize}
        \item Un paramètre scalaire est noté en italique : $K$
        \item Un vecteur de paramètres est noté en gras : $\PTh$
        \item Une matrice est notée en gras souligné : $\underline{\underline{\mathbf M}}$
        \item Un vecteur surmonté d'une flèche indique une série temporelle \textit{(généralement une température ou un flux de chaleur)} : $\TAir{t}, \PhiEm{t}$ \ldots
        \item Un symbole "chapeau" indique que la donnée est prédite par un modèle, par opposition à une donnée réelle : $\ThatAir{t}, \PhihatEm{t}$ \ldots
        \item Les fonctions usuelles et les opérateurs sont notés sans italique : $\cos, \mathrm{RMSE}$ \ldots
        \item Les unités physiques sont notées en gras : $19^\circ \mathbf C, 1000\mathbf W \ldots$
    \end{itemize}

    Ces notations sont susceptibles de varier légèrement dans la documentation du projet IBIS disponible en \ref{ch:ibisdoc}.

\end{abstract}
